Para la realización de este trabajo de fin de grado, se han revisado las principales 
plataformas de valoración y gestión de videojuegos. En general, uno de los aspectos más 
mejorables detectados en estas plataformas es el sistema de búsqueda y la accesibilidad 
a los filtros, elementos clave para una navegación eficiente.

\paragraph{RAWG}
RAWG es una plataforma web de videojuegos que permite a los usuarios descubrir títulos 
y consultarlos según su plataforma o popularidad. Además, ofrece información detallada sobre cada videojuego,
como capturas de pantalla, valoraciones de usuarios, enlaces a páginas web donde pueden adquirirse y las consolas
en las que están disponibles. También cuenta con una API que permite a los desarrolladores acceder a estos datos 
para integrarlos en sus propias aplicaciones. 

Sin embargo, RAWG resenta limitaciones en su sistema de búsqueda, ya que no permite aplicar varios filtros de forma 
simultánea. Si un usuario desea filtrar los videojuegos simultáneamente por género y por año de lanzamiento, no 
puede hacerlo de forma directa. En su lugar, debe realizar búsquedas secuenciales: primero seleccionar el género 
y, posteriormente, filtrar por año.

\paragraph{IGDB}
IGDB (Internet Game Database) es una base de datos de videojuegos lanzada por Twitch en 2014 \cite{IGDB} que almacena información 
sobre miles de títulos. Una de sus características destacadas es un contador para los videojuegos más esperados, 
que muestra los días, horas y minutos restantes para su lanzamiento. Además de permitir valorar videojuegos, IGDB ofrece 
una funcionalidad adicional: los usuarios pueden valorar los comentarios de otros mediante un sistema de 'me gusta' o 'no me gusta', 
de forma similar a los comentarios en YouTube.

En cuanto al sistema de búsqueda, presenta una mejora en este aspecto, ya que permite aplicar múltiples filtros en una sola búsqueda. 
Sin embargo, la accesibilidad a estos filtros no es óptima, pues el usuario debe primero hacer clic en la barra de búsqueda y luego seleccionar 
“Advanced search” para que se desplieguen las opciones de filtrado.

\paragraph{Backloggd}
Backloggd es una plataforma web que ofrece funcionalidades similares a las de RAWG e IGDB, como la valoración y organización de videojuegos. 
Una de sus características más destacadas es la posibilidad de responder a los comentarios de otros usuarios, lo que fomenta la interacción 
y el debate en torno a aspectos positivos y negativos de los videojuegos. Además, cuando un título forma parte de una saga o trilogía, la 
plataforma muestra los demás juegos relacionados, facilitando al usuario el seguimiento completo de la serie.

En cuanto a los filtros, Backloggd presenta un flujo de navegación más complejo. Para acceder a los filtros, el usuario debe ir 
a la sección “Games”, hacer clic en “Apply Filters” y, a continuación, seleccionar los criterios deseados, lo que añade pasos adicionales
que podrían simplificarse.

